\documentclass[12pt]{article}
\usepackage{fullpage,epsfig}
\usepackage{natbib}
\usepackage{hyperref}
\usepackage{bm,dsfont,hyperref}

\newcommand{\mb}[1]{\mathbf{#1}}
\newcommand{\ds}[1]{\mathds{#1}}

\begin{document}
\newcommand{\vect}[1]{\boldsymbol{#1}}
\renewcommand{\baselinestretch}{1.5} \small\normalsize

\noindent {\bf``Bayesian Statistical Inversion for High-Dimensional Computer Model
Output and Spatially Distributed Counts''}

\bigskip
\noindent
Dear Dr. Tuo,

\medskip Thank you for allowing us to revise our manuscript. We 
are thrilled that you think it is a good fit for \textit{IISE Transactions}. We
are grateful for the thoughtful suggestions made by each of the reviewers and have
addressed their comments in our revision.  We believe that addressing these concerns
has led to an improved manuscript and hope that you will agree.  I have included
responses to each comment made by a reviewer in the following pages. Additional text
is highlighted in blue in the revised manuscript.

Please let me know if you need anything else from us, and thank you again for your
consideration of our manuscript.

\vspace{0.25in}
\begin{center}
\begin{tabular}{l}
Sincerely,\\
\\
Steven Barnett
\end{tabular}
\end{center}

\bigskip
\renewcommand{\baselinestretch}{1.1} \small\normalsize

\pagebreak

\bigskip
\noindent
1. \textbf{Associate Editor} \textit{I agree with both referees’ comments that
the notation needs to be improved. When introducing u, \(X\), and \(Y^F\) on
page 5, these vectors/matrices could be introduced together with their
dimensions (e.g., u in \(R^p\) (or u in \(\Theta \subset R^p\)), \(X\) in
\(R^{(?)}\), \(X^F\) in \(R^{(?)}\), and \(Y^F\) in \(R^{(n_F?)}\)). If the
authors adopt a consistent convention, such as using lower-case letters for
vectors and upper-case letters for matrices and specifying their dimensions,
readers can more easily follow the notation.}

\bigskip
\noindent
TBD

\bigskip
\noindent
2. \textbf{Associate Editor} \textit{In addition, when introducing the dimensions
of vectors and matrices on page 5, I suggest explicitly relating these numbers
to your application. I understand you mentioned this later; however, this
would help emphasize to readers that the problem you are addressing is
large-scale. e.g., what is the size of \(Y^F, X^F\), and u in your example?}

\bigskip
\noindent
TBD

\bigskip
\noindent
3. \textbf{Associate Editor} \textit{The authors refer to the rows of \(X^F\);
however, it should be clarified earlier what each row and column of this
matrix represents when it is introduced early in Section 2. Perhaps this will
be resolved once the dimensions are introduced and clarified.}

\bigskip
\noindent
TBD

\bigskip
\noindent
4. \textbf{Associate Editor} \textit{Later, when describing the toy example in
Figure 3, the stars, circles, and the dimensions of \(X, u, Y^F, X^F\) could
be explicitly linked to the notation introduced earlier. For example, it would
help to clarify whether the circles or the stars correspond to the \(y_i\) in
Algorithm 1.}

\bigskip
\noindent
TBD

\bigskip
\noindent
5. \textbf{Associate Editor} \textit{Page 7, line 7: I understand that
\(u^{(t)}\) corresponds to a draw from Metropolis sampling. However, it has
not introduced at this point.}

\bigskip
\noindent
TBD

\bigskip
\noindent
6. \textbf{Associate Editor} \textit{On page 10, it is unclear what ``n'' refers
to. Does it correspond to the size of any of the previously introduced
notation (e.g., \(X, Y^F, u\)), or is it intended as a generic sample size for
computer runs? It is used three times on this page.}

\bigskip
\noindent
TBD

\bigskip
\noindent
7. \textbf{Associate Editor} \textit{Figure 5: Is it possible to provide a heat
map for the residuals between actual and predicted outputs?}

\bigskip
\noindent
TBD

\bigskip
\noindent
8. \textbf{Associate Editor} \textit{On page 15, please introduce an abbreviation
for your method in the text, e.g., Scaled Vecchia (SVEC).}

\bigskip
\noindent
TBD

\bigskip
\noindent
9. \textbf{Associate Editor} \textit{I agree with one referee’s comment regarding
potential computational bottlenecks in SEPIA due to the number of simulation
runs. It would strengthen the manuscript if the authors could discuss whether
approaches such as those proposed by Higdon et al., combined with Vecchia-type
approximations, could help alleviate these computational challenges.}

\bigskip
\noindent
TBD

\bigskip
\noindent
10. \textbf{Associate Editor} \textit{Relatedly, since log-transformed count data
are sometimes modeled using Gaussian likelihoods, it would be helpful if the
authors could comment on how the results might change under such a
transformation.}

\bigskip
\noindent
TBD

\bigskip
\noindent
\textbf{Associate Editor} \textit{I agree with one referee's comment about the typos need to be fixed. Some of them I found:}

\begin{itemize}
	\item \textit{Page 2, line 45: Use Figure 1 instead of Fig. 1}
	\item \textit{Page 2/8, line 52/19: Use i.e. not i.e}
\end{itemize}

\bigskip
\noindent
We greatly appreciate your careful attention to detail during your review.
These typos have been corrected.

\bigskip
\noindent
12. \textbf{Reviewer 1:} \textit{Section 2. The major writing consistency issues
in notations seem to arise in this section. While Eq. 1 provides a good
baseline for the standard Gaussian likelihood inverse problem, the undefined
variables and notations are odd. For examples, }

\begin{itemize}
	\item $N_{n_F} (\cdot, \cdot)$ \textit{is not defined here, but defined
much later in Section 2.2 introducing Gaussian process surrogates.}
	\item \textit{If \(\rho, V\) are unimportant to the discussion (which I believe
they are), perhaps do not state them, or rather provide a brief statement.}
	\item \textit{The dimension p seems to have never been defined.}
\end{itemize}

\bigskip
\noindent
TBD

\bigskip
\noindent
13. \textbf{Reviewer 1:} \textit{Section 2. In describing Metropolis sampling
``using a multivariate normal likelihood whose mean} $\mu(X^F) = m(u^{(t)},
X^F) \ldots,$'' $u^{(t)}$ \textit{seems to suggest the tth step sample collected,
but t was never defined before the MCMC algorithm.}

\bigskip
\noindent
TBD

\bigskip
\noindent
14. \textbf{Reviewer 1:} \textit{Section 2 (and more generally). The observational
data suggest that there are gaps in the data, i.e., data are not available at
every pixel. Do the computer models generate any incomplete output for certain
parameters? Moreover, are the missing data accounted for in the likelihood?}

\bigskip
\noindent
TBD

\bigskip
\noindent
15. \textbf{Reviewer 1:} \textit{Section 2.1 provide an educational illustration
of the difference between a Gaussian and a Poisson likelihood. Should readers
expect the posterior draws to cover the counts/exposure data? Or is it a
deficiency of the computer model in such case?}

\bigskip
\noindent
TBD

\bigskip
\noindent
16. \textbf{Reviewer 1:} \textit{Section 2.3. ``However, there is a method
tailor-made for our Bayesian inverse problem setting when the likelihood (for
\(Y^F\)) is Gaussian.''}

\bigskip
\noindent
\textit{This seems to be an important justification of the
proposed method, but missing words made the point rather unclear. Do the
authors mean ``there is not a method ... for Poisson likelihood''?}

\bigskip
\noindent
\textit{Also, should \(F\) in \(Y^F\) be in the superscript?}

\bigskip
\noindent
TBD

\bigskip
\noindent
17. \textbf{Reviewer 1:} \textit{Section 2.3. In the software survey, the authors
have mentioned SEPIA by Gattiker et al., 2020. I believe at two other possible
implementations are available from 2021 onward:}

\begin{itemize}
	\item \textit{GPflow (Van der Wilk et al., 2020), in addition to GP approximations, seems to also
offer a Poisson likelihood implementation.}
	\item \textit{surmise (Plumlee et al., 2025) contains an implementation of the Higdon et al. (2008)
strategy. Examples using the package with the relevant high-dimensional output
meth- ods are found, for examples, in Chan et al. (2024); S¨urer et al.
(2022).}
\end{itemize}

\bigskip
\noindent
\textit{And, what about any variational methods, e.g., Hensman et al. (2015)?}

\bigskip
\noindent
TBD

\bigskip
\noindent
18. \textbf{Reviewer 1:} \textit{Section 3.2. Since the surrogate outputs are
high-dimensional, would a multivariate measurement of performance be more
appropriate? RMSE and CRPS are univariate measures.}

\bigskip
\noindent
TBD

\bigskip
\noindent
19. \textbf{Reviewer 1:} \textit{e.g. Section 4.2. In several sections when the
surrogate prediction of a full sky map is compared with the computer model,
the authors state for multiple times that there are no ``visual'' difference. In
such cases, a residual plot seems to be beneficial to better quantify the
claim.}

\bigskip
\noindent
TBD

\bigskip
\noindent
20. \textbf{Reviewer 1:} \textit{Section 5. The demonstration of a need for a
discrepancy model is good, yet the value of Figure 13 is not clear to me, as
most posterior region becomes invisible.}

\bigskip
\noindent
TBD

\bigskip
\noindent
21. \textbf{Reviewer 1:} \textit{Computation. While the code repository is
relatively complete, it is helpful to include the software requirements and
dependencies somewhere, e.g., GPvecchia, tmvtnorm, laGP, R.utils, etc.}

\bigskip
\noindent
TBD

\bigskip
\noindent
23. \textbf{Reviewer 1:} \textit{Some language choices seem to leave
interpretative room with regards to the authors’ intentions when a reader
will benefit from a clear opinion from the authors. For example, the adverb
''extremely'' is used in several places where some quantification will be
helpful to answer the question ``in what way is it extreme?''}

\begin{itemize}
	\item \textit{``These simulations are extremely sophisticated and ...''
	(This usage seems fine.)}
	\item \textit{``We find the Vecchia approximation to be extremely effective
	in this regard.''}
	\item \textit{``GP surrogates are extremely effective.''}
	\item \textit{``the computer model output is extremely high-dimensional
	...'' A reminder of the number of outputs may be helpful here.}
\end{itemize}

\bigskip
\noindent
TBD

\bigskip
\noindent
24. \textbf{Reviewer 1:} \textit{The intention of quote usages is unclear at times.}

\begin{itemize}
	\item \textit{``status quo'' in correspondence to SEPIA.}
	\item \textit{``true'' parameter, as in synthetically generated true values?}
\end{itemize}

\bigskip
\noindent
TBD

\bigskip
\noindent
25. \textbf{Reviewer 1:} \textit{Section 1. The use of nouns with different
meanings, i.e., ``challenges'', ``bottlenecks'', and ``situations'', to refer to the
same set of issues has made it difficult to follow.}

\bigskip
\noindent
TBD

\bigskip
\noindent
26. \textbf{Reviewer 1:} \textit{Would combining Figures 2 and 3 as a top and
bottom panel help with the illustration? The figures are similar and would
benefit from side-by-side comparison.}

\bigskip
\noindent
TBD

\bigskip
\noindent
27. \textbf{Reviewer 1:} \textit{The legend in Figure 7 (right panel) is unclear.}

\bigskip
\noindent
TBD

\bigskip
\noindent
28. \textbf{Reviewer 1:} \textit{The caption of Figure 8 can be improved,
especially regarding the laGP line in the rightmost panel.}

\bigskip
\noindent
TBD

\bigskip
\noindent
\textbf{Reviewer 1:} \textit{Supplemental comments}

\bigskip
\noindent
Thank you for your diligent reading of our manuscript. We have fixed the typos
and phrasing issues mentioned. We also now introduce the KOH acronym earlier
in the paper.

\bigskip
\noindent
30. \textbf{Reviewer 2:} \textit{The manuscript uses \(X^M\) to denote the
combined input \((u, X)\) to the surrogate, while \(X^F\) denotes the
field-data locations (latitude/longitude) and does not involve \(u\). Since
\(M\) indicates ``model,'' I found this overloading of \(X\) (and the use of
superscripts) easy to misread: I repeatedly had to remind myself that \(X^F\)
is purely spatial, whereas \(X^M\) includes both parameters and locations. I
recommend reserving \(X^M\) (or simply \(X\)) for model output locations only,
and introducing a different symbol for the joint design in parameter–location
space (e.g., \(Z=(u,X)\)). This small change would reduce cognitive load and
make the distinction between ``where data are observed'' and ``what the emulator
is conditioned on'' more transparent.}

\bigskip
\noindent
TBD

\bigskip
\noindent
31. \textbf{Reviewer 2:} \textit{Figure 3 is helpful overall, but I suggest
removing the star markers in the middle panel, since they do not correspond to
the final quantity of interest. It would make the visual assesment of the
inference performance easier.}

\bigskip
\noindent
TBD

\bigskip
\noindent
32. \textbf{Reviewer 2:} \textit{In Figure 5, it would be helpful to annotate each
simulated output and its surrogate approximation with the corresponding
parameter value(s) used to generate them (either in the panel titles, caption,
or a small table/legend). This would make the comparison more interpretable.}

\bigskip
\noindent
TBD

\bigskip
\noindent
33. \textbf{Reviewer 2:} \textit{The calibration parameters are described
essentially in terms of \((u_1, u_2/u_1)\). Could the authors clarify the
motivation for using this parameterization, as opposed to solving the inverse
problem directly in terms of \((u_1, u_2)\) with \(u_1\) and \(u_2\) treated
as independent coordinates? I understand there is physical meaning of the
ratio parameter. But isn’t it a general desire to use untangled variables in
regression problem?}

\bigskip
\noindent
TBD

\bigskip
\noindent
34. \textbf{Reviewer 2:} \textit{Figure 6 suggests that SEPIA (even with only
three PC bases) can achieve performance comparable to SVEC (this acronymm is
never defined in the text). If the primary bottleneck for SEPIA in this
setting is scalability with a large number of simulation runs (Figure 8,
right panel), a potentially manageable alternative would be to retain the
Higdon et al.basis-decomposition setup but model the PC/basis weights using
Vecchia-approximated GPs (rather than full GPs). It would be helpful if the
authors could comment on whether such a ``Higdon et al. + Vecchia'' hybrid has
been considered, and how it might compare to SVEC in terms of computational
cost and accuracy. This could potentially be implemented without relying on
the SEPIA library (thereby avoiding the MCMC cost of full Bayesian inference),
and instead used as a drop-in surrogate within the proposed framework.}

\bigskip
\noindent
TBD

\bigskip
\noindent
35. \textbf{Reviewer 2:} \textit{Since it is common to model log-count data using
a Gaussian likelihood, and the IBEX application uses only 66 simulator runs,
it would be helpful if the authors could comment on, and possibly experiment
whether SEPIA would be computationally feasible at this problem size when
applied to (transformed) log-count outputs. Is the main limitation of SEPIA
here the non-Gaussian (Poisson) observation model?}

\bigskip
\noindent
TBD

\bigskip
\noindent
36. \textbf{Reviewer 2:} \textit{It remains unclear to me how the ``stacked'' field
data are used in the real-data analysis (e.g., stacking three years into
roughly 23K rows of \((X^F, Y^F)\)). Figure 11 suggests that \(X^F\) varies
across years. Does the simulator also produce year-specific outputs (at
year-specific \(X^F\)), which are then stacked and passed into a single
likelihood evaluation, or are the years treated as separate likelihood terms
(with a product over years)? A short mathematical description (with explicit
notation for year index, locations, and simulator output) would be very
helpful for readers.}

\bigskip
\noindent
TBD

\end{document}
